%% 6-11 双曲線関数 cosh・sinh（第14回：組 V6-6。ページ制限なし。偶数ページ）
%% ほかの巻が指している形：v = v_t tanh(gt/v_t)、tanh の読み方「タンジェント・エイチ」、(tanh u)' = 1 − tanh² u（物理 力-2）、
%%   スカイダイバー 80 kg・c = 0.26・v_t ≈ 54.9 m/s（9.70・18.8・26.9・39.1・51.9）、6-3 のつり橋（放物線）と対になる懸垂線
%% 3段：目でわかる（くさりの形・e^x と e^{-x} の平均と差）→手でわかる（定義・cosh²−sinh²=1）→式でわかる（微分・y''=a²y・落下・懸垂線）
\makeatletter
\newcommand\bkvilhead[1]{\par\addvspace{6pt}\Needspace{6\baselineskip}\noindent{\sffamily\bfseries\color{\bkTT}#1}\par\nobreak\vspace{2pt}}
\newcommand\bkvilbox[1]{\par\addvspace{3pt}\noindent\textbf{#1}\par\noindent\ignorespaces}
\newcommand\bkvilframe[1]{\par\addvspace{4pt}{\centering\tikz\node[draw=\bkTT, line width=0.9pt, fill=bkTheme!10, rounded corners=1.5mm, inner xsep=6mm, inner ysep=2.5mm, align=center]{#1};\par}\addvspace{4pt}\noindent\ignorespaces}
\newcommand\bkvilline[1]{\par\addvspace{3pt}{\centering #1\par}\addvspace{3pt}\noindent\ignorespaces}
\newcommand\bkvilq{\par\vspace{2pt}\noindent\hangindent=4.3\zw\hangafter=1\hspace*{2.6\zw}}
\newcommand\bkvila{\par\vspace{2pt}\noindent\hangindent=2.6\zw\hangafter=0\relax}
\newcommand\bkvilb{\par\vspace{2pt}\noindent\hangindent=4.3\zw\hangafter=0\relax}

% 電柱  \bkvilpole{(x,y)}{高さ}
\newcommand\bkvilpole[2]{\begin{scope}[shift={#1}]
  \fill[bkGray!60] (-0.9,0) rectangle (0.9,#2);
  \draw[bkGray!80!black, line width=0.6pt] (-0.9,0) rectangle (0.9,#2);
  \fill[bkGray!80!black] (-3.2,{#2-2.2}) rectangle (3.2,{#2-1.4});
\end{scope}}
\makeatother

\bkchapter{6}{11}{双曲線関数 cosh・sinh}{\bklevelrange{4}{5}}{双曲線関数 cosh・sinh}
\bkmichishirube{5-14,5-17,5-18,6-3,6-5,6-10}{6-12}

\begin{goal}
\begin{itemize}[label={\color{\bkTT}●}, leftmargin=1.6\zw, itemsep=0pt, topsep=1pt]
\item \bklevel{4}\enspace $\cosh x$・$\sinh x$・$\tanh x$ を $e^x$ と $e^{-x}$ で表し、値を求めて、グラフの形を説明できる。
\item \bklevel{4}\enspace $\cosh^2x-\sinh^2x=1$ と、$\cosh$・$\sinh$・$\tanh$ の微分のきまりを使える。
\item \bklevel{5}\enspace $y''=a^2y$ の答えを $\cosh$・$\sinh$ で書き、空気の中の落下や、たるんだ電線の形に使える。
\end{itemize}
\end{goal}

\begin{hatena}
\noindent\begin{minipage}[c]{0.6\linewidth}
\raggedright
電柱と電柱のあいだにかけた電線や、両手でつるしたネックレスは、なめらかにたるみます。その形は、3-9 で学んだ放物線にそっくりです。ガリレオも、1638年の本で、つるしたくさりの形は放物線に近いと書きました。

ところが、くわしく調べると、この形は放物線ではありません。いったい、何という曲線なのでしょう。
\end{minipage}\hfill
\begin{minipage}[c]{0.37\linewidth}
\centering
\begin{tikzpicture}[x=0.9mm, y=1mm]
  \draw[bkGray, line width=0.8pt] (-2,0) -- (48,0);
  \bkvilpole{(2,0)}{34}
  \bkvilpole{(44,0)}{34}
  \draw[bkOrange, line width=0.9pt, dash pattern=on 2.5pt off 1.5pt, domain=-21:21, samples=60] plot ({23+\x}, {17+15*(\x/21)^2});
  \draw[bkInk, line width=1.3pt, domain=-21:21, samples=80] plot ({23+\x}, {17+16.733*(cosh(\x/16.733)-1)});
  \node[font=\sffamily\bfseries\Large, text=bkYellow!40!black] at (23,27) {？};
\end{tikzpicture}
\end{minipage}
\end{hatena}

\noindent 6-10 で、$\cos\theta=\dfrac{e^{j\theta}+e^{-j\theta}}{2}$ という式を作りました。この章では、この式から $j$ を外した形 $\dfrac{e^{x}+e^{-x}}{2}$ を調べます。この関数が、たるんだくさりの形をえがき、空気の中を落ちるスカイダイバーの速さ（6-5）を表し、5-18 の「たおれる棒」の式もすっきり書けることを見ていきます。

この章も大学の内容です。「目でわかる」「手でわかる」の段だけでも、新しい関数の形と使い方がわかります。1回で分からなくても大丈夫です。

%% ---------------- 思い出そう・やってみよう ----------------
\begin{omoidasou}
\bkrecall{3-9}{放物線 $y=ax^2$。}
\bkrecall{5-14}{$(e^{ax})'=ae^{ax}$。$e^{-x}=\dfrac{1}{e^x}$。$e^0=1$。}
\bkrecall{5-17}{$e^x=1+x+\dfrac{x^2}{2!}+\dfrac{x^3}{3!}+\cdots$}
\bkrecall{5-18}{$f''=a^2f$ の答えは $Ae^{at}+Be^{-at}$。たおれる棒 $\theta=0.005(e^{2t}+e^{-2t})$。}
\bkrecall{6-3}{つり橋のケーブルは、$\dfrac{\dd^2y}{\dd x^2}$ が一定なので放物線。両はしの状態（境界条件）で数を決める。}
\bkrecall{6-4}{変数分離：$y$ をふくむ項を左に、$t$ を右に集めて、両辺を積分する。}
\bkrecall{6-5}{速さの2乗に比例する抵抗：$m\dfrac{\dd v}{\dd t}=mg-cv^2$。終端速度 $v_{\mathrm t}=\sqrt{\dfrac{mg}{c}}$。\\スカイダイバー（\qty{80}{kg}、$c=\qty{0.26}{N.s^2/m^2}$）で $v_{\mathrm t}\approx\qty{54.9}{m/s}$。}
\bkrecall{6-10}{$\cos\theta=\dfrac{e^{j\theta}+e^{-j\theta}}{2}$、$\sin\theta=\dfrac{e^{j\theta}-e^{-j\theta}}{2j}$。}
\end{omoidasou}

\begin{bkfigure}
\begin{tikzpicture}[x=1mm, y=1mm, nd/.style={bk node, font=\sffamily\footnotesize, text width=25mm, inner sep=1.5mm, minimum height=10mm}]
  \node[nd] (a) at (0,0) {① $e^x$ と $e^{-x}$ の\\平均と差の半分};
  \node[nd, draw=\bkTT, line width=1.2pt, fill=bkTheme!12] (b) at (34,0) {② $\cosh$・$\sinh$\\・$\tanh$};
  \node[nd] (c) at (68,0) {③ 2乗の差が1\\（双曲線）};
  \node[nd] (d) at (102,0) {④ 微分の\\きまり};
  \node[nd] (e) at (34,-17) {⑤ $y''=a^2y$\\を解く};
  \node[nd] (f) at (68,-17) {⑥ 空気の中の\\落下（$\tanh$）};
  \node[nd] (g) at (102,-17) {⑦ くさりの形\\（懸垂線）};
  \draw[bk link] (a) -- (b);
  \draw[bk link] (b) -- (c);
  \draw[bk link] (c) -- (d);
  \draw[bk link] (d.south) -- ++(0,-3.5) -| (e.north);
  \draw[bk link] (e) -- (f);
  \draw[bk link] (f) -- (g);
\end{tikzpicture}
\bkfigcaption{この章の道すじ。$e^x$ と $e^{-x}$ から新しい関数を作り、落下とくさりの形の2つの問題を解く。}
\end{bkfigure}

\begin{yattemiyou}[ネックレスをつるして、形を写しとろう]
\textbf{用意するもの}\par
細いくさりのネックレス、方眼紙、コルクボード、がびょう2本
\begin{enumerate}[label={\bkcircled[bkGreen]{\arabic*}}, leftmargin=2.4\zw, itemsep=1pt, topsep=2pt]
\item 方眼紙をコルクボードにはり、同じ高さに、20\,cm はなしてがびょうを2本さす。
\item ネックレスをがびょうにかけて、いちばん下が、がびょうより 15\,cm 下になるようにたらす。
\item いちばん下の点を原点にして、横に 2\,cm ごとに、くさりの高さを読む。
\item 同じはばとたるみの放物線 $y=0.15x^2$（$x=10$ で15）の値とくらべる。
\end{enumerate}
\noindent\begin{minipage}[c]{0.5\linewidth}
\centering
\begin{bktable}{colspec={Q[c,m,wd=17mm] Q[c,m,wd=14mm] Q[c,m,wd=14mm]}, rowsep=1.5pt, colsep=3pt}
横 $x$（cm） & くさり（cm） & 放物線（cm） \\
0 & 0 & 0 \\
2 & 0.4 & 0.6 \\
4 & 1.8 & 2.4 \\
6 & 4.3 & 5.4 \\
8 & 8.5 & 9.6 \\
10 & 15.0 & 15.0 \\
\end{bktable}
\bktabcaption{ソラさんがはかった結果}
\end{minipage}\hfill
\begin{minipage}[c]{0.48\linewidth}
\centering
\begin{tikzpicture}[x=2.5mm, y=2.5mm]
  \foreach \i in {-10,-8,...,10} {\draw[bkLine, line width=0.3pt] (\i,0) -- (\i,16);}
  \foreach \k in {0,2,...,16} {\draw[bkLine, line width=0.3pt] (-10,\k) -- (10,\k);}
  \draw[bkOrange, line width=1pt, dash pattern=on 2.5pt off 1.5pt, domain=-10:10, samples=60] plot (\x, {0.15*\x*\x});
  \draw[bkInk, line width=1.8pt, line cap=round, dash pattern=on 0pt off 2.2pt, domain=-10:10, samples=120] plot (\x, {4.7568*(cosh(\x/4.7568)-1)});
  \fill[bkRed] (-10,15) circle (0.35); \fill[bkRed] (10,15) circle (0.35);
  \node[font=\sffamily\footnotesize, anchor=north, inner sep=1pt] at (0,-0.4) {原点};
  \node[font=\sffamily\footnotesize, text=bkOrange!80!black, anchor=north west, inner sep=1pt, fill=white] at (-7.5,15.6) {放物線};
  \draw[bkInk, line width=0.4pt] (6.5,3.65) -- (4.8,8.6) node[anchor=south, font=\sffamily\footnotesize, inner sep=1pt, fill=white] {くさり};
\end{tikzpicture}
\end{minipage}
\end{yattemiyou}

\noindent はしと、いちばん下の点はそろえたのに、とちゅうでは、くさりのほうが放物線より最大1\,cm ほど低くなりました。くさりの形は、底のあたりが平らで、はしで急に立ち上がります。この形を式で表すのが、この章の目標の1つです。



%% ---------------- 数学の道具 ----------------
\bksection{doguu}{数学の道具}

\bkstage{me}{目でわかる}{$e^x$ と $e^{-x}$ を組み合わせる}{\bklevel{4}}

\bksubsection{1}{ふえる $e^x$ と、へる $e^{-x}$}

\label{bk:6-11:heikin}$e^x$ は右へ行くほどふえ、$e^{-x}$ は右へ行くほどへります。$e^{-x}$ のグラフは、$e^x$ のグラフを、たての軸で折り返したものです（$x$ を $-x$ にかえると、左右が入れかわる）。この2つから、新しい関数を2つ作ります。
\begin{itemize}[leftmargin=1.2\zw, itemsep=2pt, topsep=2pt, label={\color{\bkTT}●}]
\item \textbf{平均}\quad $\dfrac{e^{x}+e^{-x}}{2}$（2つの高さの、ちょうどまん中）
\item \textbf{差の半分}\quad $\dfrac{e^{x}-e^{-x}}{2}$
\end{itemize}
表で計算してみます（$e^{1}=2.718$、$e^{-1}=0.368$ など）。
\begin{bkcenter}
\begin{bktable}{colspec={Q[l,m,wd=30mm] *{5}{Q[c,m,wd=14mm]}}, rowsep=1.5pt, colsep=2pt}
$x$ & $-2$ & $-1$ & 0 & 1 & 2 \\
$e^x$ & 0.135 & 0.368 & 1 & 2.718 & 7.389 \\
$e^{-x}$ & 7.389 & 2.718 & 1 & 0.368 & 0.135 \\
平均 $\dfrac{e^{x}+e^{-x}}{2}$ & 3.762 & 1.543 & 1 & 1.543 & 3.762 \\
差の半分 $\dfrac{e^{x}-e^{-x}}{2}$ & $-3.627$ & $-1.175$ & 0 & 1.175 & 3.627 \\
\end{bktable}
\bktabcaption{平均は左右で同じ値。差の半分は、左右で符号だけが反対。}
\end{bkcenter}

\begin{bkfigure}
\begin{tikzpicture}
\begin{axis}[bkplot, scale only axis, width=92mm, height=58mm, xmin=-2.6, xmax=2.6, ymin=-6.2, ymax=6.4, xtick={-2,-1,1,2}, ytick={-4,-2,2,4,6}, xlabel={$x$}, ylabel={$y$},
    yticklabel style={fill=white, inner sep=1pt, outer sep=2.3pt}, xticklabel style={fill=white, inner sep=1pt}]
  \addplot[bkGray, line width=0.9pt, domain=-2.6:1.82, samples=60] {exp(x)};
  \addplot[bkGray, line width=0.9pt, dashed, domain=-1.82:2.6, samples=60] {exp(-x)};
  \addplot[bkcurve, line width=2pt, domain=-2.38:2.38, samples=80] {(exp(x)+exp(-x))/2};
  \addplot[bkcurve2, line width=2pt, domain=-2.5:2.5, samples=80] {(exp(x)-exp(-x))/2};
  \node[font=\sffamily\footnotesize, text=bkGray, anchor=west, inner sep=1pt, fill=white] at (axis cs:1.86,6.0) {$e^x$};
  \node[font=\sffamily\footnotesize, text=bkGray, anchor=east, inner sep=1pt, fill=white] at (axis cs:-1.86,6.0) {$e^{-x}$};
  \node[font=\sffamily\footnotesize, text=bkBlue, anchor=west, inner sep=1pt, fill=white] at (axis cs:0.75,2.75) {平均};
  \node[font=\sffamily\footnotesize, text=bkOrange!80!black, anchor=west, inner sep=1pt, fill=white] at (axis cs:1.05,0.55) {差の半分};
\end{axis}
\end{tikzpicture}
\bkfigcaption{$e^x$ と $e^{-x}$（灰色）の平均（青）と差の半分（オレンジ）。平均は谷の形、差の半分は原点を通ってふえ続ける。}
\end{bkfigure}

\bkvilhead{図から読みとれること}
\begin{enumerate}[label={\bkcircled[\bkTF]{\arabic*}}, leftmargin=2.4\zw, itemsep=2pt, topsep=2pt]
\item \textbf{平均は、いちばん低い所が1}\quad $x=0$ では、どちらも $e^0=1$ なので平均は1。$x$ が0からはなれると、片方がどんどん大きくなり、もう片方は小さくなっても0より下にはならないので、平均は1より大きくなります。左右対称の谷の形で、くさりの形に似ています。
\item \textbf{差の半分は、原点を通る}\quad $x=0$ で $1-1=0$。右では $e^x$ のほうが大きいので正、左では負。原点について対称で、ずっとふえ続けます。
\item \textbf{遠くでは、2つがくっつく}\quad $x$ が大きいと $e^{-x}$ はほとんど0なので、平均も差の半分も、ほとんど $\dfrac{e^x}{2}$ になります。たとえば $x=3$ では、平均は10.068、差の半分は10.018 です。
\end{enumerate}

%% ---------------- 手でわかる ----------------
\Needspace{16\baselineskip}
\bkstage{te}{手でわかる}{名前をつけて、計算する}{\bklevel{4}}

\bksubsection{2}{双曲線関数 $\cosh$・$\sinh$・$\tanh$}

\label{bk:6-11:teigi}①の2つの関数と、その比に、名前をつけます。
\bkvilframe{$\cosh x=\dfrac{e^{x}+e^{-x}}{2}$\qquad $\sinh x=\dfrac{e^{x}-e^{-x}}{2}$\qquad $\tanh x=\dfrac{\sinh x}{\cosh x}$}
まとめて\textbf{双曲線関数}といいます。\bkidx{そうきょくせんかんすう}{双曲線関数}「双曲線」という名前のわけは、③で分かります。読み方は次のとおりです。
\begin{bkcenter}
\begin{bktable}{colspec={Q[c,m,wd=16mm] Q[l,m,wd=50mm] Q[l,m,wd=40mm]}, rowsep=1.5pt}
記号 & 読み方 & 6-10 とのちがい \\
$\cosh x$ & ハイパボリック・コサイン & $\cos\theta=\dfrac{e^{j\theta}+e^{-j\theta}}{2}$ \\
$\sinh x$ & ハイパボリック・サイン & $\sin\theta=\dfrac{e^{j\theta}-e^{-j\theta}}{2j}$ \\
$\tanh x$ & ハイパボリック・タンジェント（タンジェント・エイチとも読む） & $\tan\theta=\dfrac{\sin\theta}{\cos\theta}$ \\
\end{bktable}
\bktabcaption{「ハイパボリック（hyperbolic）」は「双曲線の」という意味。$\cos$ の後ろの h は、その頭文字。}
\end{bkcenter}
$\cos\theta$ の式から $j$ をとると $\cosh x$ になります。$\sin\theta$ の式は、分母の $2j$ の $j$ もとると $\sinh x$ になります。三角関数の「いとこ」のような関数です。

\bkvilhead{値を求める}
表計算では \texttt{=COSH(A2)}、\texttt{=SINH(A2)}、\texttt{=TANH(A2)} で計算できます。関数電卓では、hyp のキーを押してから cos・sin・tan のキーを押します。
\begin{bkcenter}
\begin{bktable}{colspec={Q[c,m,wd=16mm] *{6}{Q[c,m,wd=13.5mm]}}, rowsep=1.5pt, colsep=2pt}
$x$ & 0 & 0.5 & 1 & 1.5 & 2 & 3 \\
$\cosh x$ & 1 & 1.128 & 1.543 & 2.352 & 3.762 & 10.068 \\
$\sinh x$ & 0 & 0.521 & 1.175 & 2.129 & 3.627 & 10.018 \\
$\tanh x$ & 0 & 0.462 & 0.762 & 0.905 & 0.964 & 0.995 \\
\end{bktable}
\bktabcaption{$\tanh x$ は、$x$ が大きくなると1に近づく。}
\end{bkcenter}

\bkvilhead{$\tanh x$ のグラフ}
\bkidx{はいぱぼりっくたんじぇんと}{ハイパボリック・タンジェント（tanh）}
\noindent\begin{minipage}[c]{0.5\linewidth}
\raggedright
$\tanh x=\dfrac{\sinh x}{\cosh x}$ は、分子と分母の両方が①の図の曲線です。$x$ が大きいと、2つはほとんど同じ $\dfrac{e^x}{2}$ になるので、比は1に近づきます。$x$ が大きな負の数なら、$\sinh x$ は $-\dfrac{e^{-x}}{2}$ に近づくので、比は $-1$ に近づきます。$\tanh x$ は、\textbf{いつも $-1$ と1のあいだ}にあり、原点の近くでは、$y=x$ とほとんど同じ（傾き1）です。どんな大きな数でも $-1$ から1のあいだにおさめるので、AI（ニューラルネットワーク）の中でも使われてきました。
\end{minipage}\hfill
\begin{minipage}[c]{0.46\linewidth}
\centering
\begin{tikzpicture}
\begin{axis}[bkplot, scale only axis, width=48mm, height=34mm, xmin=-3.6, xmax=3.6, ymin=-1.45, ymax=1.5, xtick={-3,-2,-1,1,2,3}, ytick={-1,1}, xlabel={$x$},
    yticklabel style={fill=white, inner sep=1pt, outer sep=2.3pt}, xticklabel style={fill=white, inner sep=1pt}]
  \addplot[bkGray, dashed, line width=0.7pt, domain=-3.6:3.6] {1};
  \addplot[bkGray, dashed, line width=0.7pt, domain=-3.6:3.6] {-1};
  \addplot[bkthin, domain=-1.3:1.3] {x};
  \addplot[bkPurple, line width=2pt, domain=-3.5:3.5, samples=80] {tanh(x)};
  \node[font=\sffamily\footnotesize, text=bkPurple, anchor=north west, inner sep=1pt, fill=white] at (axis cs:1.2,0.72) {$\tanh x$};
  \node[font=\sffamily\footnotesize, text=bkGray, anchor=east, inner sep=1pt, fill=white] at (axis cs:-0.15,1.2) {$y=x$};
\end{axis}
\end{tikzpicture}
\end{minipage}

\bksubsection{3}{$\cosh^2x-\sinh^2x=1$：名前のわけ}

\label{bk:6-11:sokyoku}三角関数には $\cos^2\theta+\sin^2\theta=1$（4-8）がありました。双曲線関数ではどうでしょう。まず数で確かめます。$x=1$ なら、
\bkvilline{$\cosh^21-\sinh^21=1.543^2-1.175^2=2.381-1.381=1.000$}
$x=2$ でも $3.7622^2-3.6269^2=14.154-13.154=1.000$。どうやら、2乗の「差」がいつも1になりそうです。

\bkvilhead{式で確かめる}
$e^x\times e^{-x}=e^0=1$ を使います（3-11 の指数法則）。
\begin{bkhenkei}
\cosh^2x=\frac{(e^x+e^{-x})^2}{4} & 定義を2乗する \\
=\frac{e^{2x}+2e^xe^{-x}+e^{-2x}}{4} & $(a+b)^2$ を展開（3-8） \\
=\frac{e^{2x}+2+e^{-2x}}{4} & $e^xe^{-x}=1$ \\
\end{bkhenkei}
同じように、$\sinh^2x=\dfrac{e^{2x}-2+e^{-2x}}{4}$（$(a-b)^2$ なので、まん中が $-2$）。ひくと、$e^{2x}$ と $e^{-2x}$ は消えて、
\bkvilframe{$\cosh^2x-\sinh^2x=\dfrac{2-(-2)}{4}=1$}

\bkvilhead{点 $(\cosh t,\ \sinh t)$ は、双曲線の上にある}
\noindent\begin{minipage}[c]{0.5\linewidth}
\raggedright
4-8 では、点 $(\cos t,\ \sin t)$ が、$\cos^2t+\sin^2t=1$ から、円 $x^2+y^2=1$ の上を動きました。

同じように、点 $(\cosh t,\ \sinh t)$ は、$\cosh^2t-\sinh^2t=1$ から、曲線 $x^2-y^2=1$ の上を動きます。この曲線を\textbf{双曲線}といいます\bkidx{そうきょくせん}{双曲線}（2本の曲線が向かい合う形の、右の1本）。これが「双曲線関数」という名前のわけです。

$t$ がふえると、点は右の曲線を下から上へ進み、ななめの点線（$y=x$、$y=-x$）に近づいていきます。
\end{minipage}\hfill
\begin{minipage}[c]{0.46\linewidth}
\centering
\begin{tikzpicture}[x=10mm, y=10mm]
  \draw[bkGray, line width=0.6pt, -{Stealth[length=3.5pt]}] (-1.6,0) -- (3.4,0) node[anchor=north east, font=\sffamily\footnotesize, inner sep=1pt, yshift=-1.2mm] {$x$};
  \draw[bkGray, line width=0.6pt, -{Stealth[length=3.5pt]}] (0,-3.1) -- (0,3.2) node[anchor=north west, font=\sffamily\footnotesize, inner sep=1pt, xshift=1.2mm] {$y$};
  \draw[bkGray, dashed, line width=0.6pt] (-0.4,-0.4) -- (3.0,3.0);
  \draw[bkGray, dashed, line width=0.6pt] (-0.4,0.4) -- (3.0,-3.0);
  \draw[bkBlue!60, line width=0.9pt] (0,0) circle (1);
  \draw[bkPurple, line width=1.7pt, domain=-1.8:1.8, samples=80] plot ({cosh(\x)}, {sinh(\x)});
  \foreach \t in {-1.5,-1,-0.5,0,0.5,1,1.5} {\fill[bkPurple] ({cosh(\t)},{sinh(\t)}) circle (0.07);}
  \node[font=\sffamily\footnotesize, anchor=west, inner sep=1pt] at ({cosh(1)+0.12},{sinh(1)-0.1}) {$t=1$};
  \node[font=\sffamily\footnotesize, anchor=west, inner sep=1pt] at ({cosh(0)+0.12},-0.22) {$t=0$};
  \node[font=\sffamily\footnotesize, anchor=west, inner sep=1pt] at ({cosh(-1)+0.12},{sinh(-1)+0.1}) {$t=-1$};
  \node[font=\sffamily\footnotesize, text=bkPurple, anchor=west, inner sep=1pt, fill=white] at (1.75,2.75) {$x^2-y^2=1$};
  \node[font=\sffamily\footnotesize, text=bkBlue, anchor=south east, inner sep=1pt, fill=white] at (-0.72,0.72) {円};
\end{tikzpicture}
\end{minipage}

\noindent 3-7 の反比例のグラフ $y=\dfrac{a}{x}$ も、双曲線の仲間です。$x^2-y^2=1$ を45°回すと、反比例のグラフと同じ形（$y=\dfrac{1}{2x}$）になります。

\bkkokomade[$\cosh$ は平均、$\sinh$ は差の半分。$\cosh^2-\sinh^2=1$。]

%% ---------------- 式でわかる ----------------
\Needspace{16\baselineskip}
\bkstage{shiki}{式でわかる}{微分と、微分方程式への応用}{\bklevel{5}}

\bksubsection{4}{微分：$\cos$・$\sin$ と、符号だけがちがう}

\label{bk:6-11:bibun}5-14 の $(e^{x})'=e^{x}$、$(e^{-x})'=-e^{-x}$ を使って、1つずつ微分します。
\begin{bkhenkei}
(\cosh x)'=\frac{(e^x)'+(e^{-x})'}{2} & 定義を項ごとに微分（定数倍・たし算） \\
=\frac{e^x-e^{-x}}{2}=\sinh x & $(e^{-x})'=-e^{-x}$ \\
\end{bkhenkei}
\begin{bkhenkei}
(\sinh x)'=\frac{(e^x)'-(e^{-x})'}{2} & 定義を項ごとに微分 \\
=\frac{e^x+e^{-x}}{2}=\cosh x & $-(-e^{-x})=+e^{-x}$ \\
\end{bkhenkei}
\begin{bkcenter}
\begin{bktable}{colspec={Q[l,m,wd=30mm] Q[c,m,wd=36mm] Q[c,m,wd=36mm]}, rowsep=1.5pt}
 & 三角関数（5-15） & 双曲線関数 \\
$\cos$ の微分 & $(\cos x)'=-\sin x$ & $(\cosh x)'=\sinh x$ \\
$\sin$ の微分 & $(\sin x)'=\cos x$ & $(\sinh x)'=\cosh x$ \\
2回微分 & 自分の $-1$ 倍（ゆれる） & 自分の $+1$ 倍（ふえる） \\
\end{bktable}
\bktabcaption{$\cosh$ の微分には、$-$ がつかない。}
\end{bkcenter}
$\cosh x$ は、2回微分するともとにもどります（$\cosh\to\sinh\to\cosh$）。$\cos x$ は、2回微分すると $-\cos x$ でした。符号が1つちがうだけで、「ゆれる」と「ふえる」に分かれます。グラフでも確かめられます。$\cosh x$ は $x>0$ でふえているので、傾き $\sinh x$ は正です。

$\cosh ax$ のように中に $a$ がかかっているときは、合成関数の微分（5-6）で、外の微分に中の微分 $a$ をかけます。
\bkvilline{$(\cosh ax)'=a\sinh ax$\qquad $(\sinh ax)'=a\cosh ax$}

\bkvilhead{$\tanh x$ の微分}
$\tanh x=\dfrac{\sinh x}{\cosh x}$ を、商の微分 $\Bigl(\dfrac{f}{g}\Bigr)'=\dfrac{f'g-fg'}{g^2}$（5-6 のもっと知りたい）で微分します。
\begin{bkhenkei}[0.58]
(\tanh x)'=\Bigl(\frac{\sinh x}{\cosh x}\Bigr)' & 定義 \\
=\frac{(\sinh x)'\cosh x-\sinh x\,(\cosh x)'}{\cosh^2x} & 商の微分 \\
=\frac{\cosh x\cosh x-\sinh x\sinh x}{\cosh^2x} & ④の2つの微分を入れる \\
=\frac{\cosh^2x-\sinh^2x}{\cosh^2x}=\frac{1}{\cosh^2x} & ③の $\cosh^2x-\sinh^2x=1$ \\
\end{bkhenkei}
また、2行めの分数を2つに分けると、$\dfrac{\cosh^2x}{\cosh^2x}-\dfrac{\sinh^2x}{\cosh^2x}=1-\tanh^2x$ とも書けます。
\bkvilframe{$(\tanh x)'=1-\tanh^2x=\dfrac{1}{\cosh^2x}$}
$x=0$ では $\tanh0=0$ なので傾きは1、$x$ が大きいと $\tanh x\approx1$ なので傾きは0に近づきます。グラフ（②）の形と合っています。「微分が、自分自身 $\tanh x$ で書ける」ことが、⑥の落下の計算で役に立ちます。

\bksubsection{5}{$y''=a^2y$ を、$\cosh$ と $\sinh$ で解く}

\label{bk:6-11:kaku}5-18 では、「2回微分すると自分の $a^2$ 倍」（$y''=a^2y$）の答えを $Ae^{at}+Be^{-at}$ と書きました。$\cosh at$ と $\sinh at$ も、2回微分すると $a^2$ 倍になります（④）。
\bkvilline{$(\cosh at)''=a^2\cosh at$\qquad $(\sinh at)''=a^2\sinh at$}
だから、線形（5-18）の考えで、たしたものも答えです。
\bkvilframe{$y''=a^2y$ の答え：\quad $y=C\cosh at+D\sinh at$}

\bkvilhead{はじめの状態から、すぐに数が決まる}
この書き方のよいところは、$t=0$ での値です。$\cosh0=1$、$\sinh0=0$ なので、
\begin{enumerate}[label={\bkcircled[\bkTF]{\arabic*}}, leftmargin=2.4\zw, itemsep=2pt, topsep=2pt]
\item はじめの値：$y(0)=C\times1+D\times0=C$。だから $C=y(0)$。
\item はじめの傾き：微分すると $y'=aC\sinh at+aD\cosh at$。$t=0$ を入れると、
\[ y'(0)=aC\times0+aD\times1=aD \]
だから $D=\dfrac{y'(0)}{a}$。
\end{enumerate}
連立方程式を解かなくても、$C$ と $D$ が1つずつ決まります。$e^{at}$ と $e^{-at}$ では、どちらも $t=0$ で1になるので、2つの式を連立して解く必要がありました（5-18）。

\bkvilhead{5-18 のたおれる棒を書きなおす}
5-18 では、次の式と、はじめの状態を考えました。
\bkvilline{$\ddot\theta=4\theta$、\quad $\theta(0)=0.01$、\quad $\dot\theta(0)=0$}
$a^2=4$ なので $a=2$。①②から、$C=\theta(0)=0.01$、$D=\dfrac{\dot\theta(0)}{2}=0$ です。だから、
\bkvilline{$\theta=0.01\cosh2t$}
5-18 の答え $0.005(e^{2t}+e^{-2t})$ と、定義から同じものです（$0.005\times2=0.01$）。

\bkvilhead{2つの書き方の行き来}
$y=C\cosh at+D\sinh at$ に定義を入れて、$e^{at}$ と $e^{-at}$ ごとにまとめると、
\bkvilline{$y=\dfrac{C+D}{2}\,e^{at}+\dfrac{C-D}{2}\,e^{-at}$}
逆に $y=Ae^{at}+Be^{-at}$ なら、$A=\dfrac{C+D}{2}$、$B=\dfrac{C-D}{2}$ をたしたりひいたりして、$C=A+B$、$D=A-B$ です。

\Needspace{20\baselineskip}
\begin{motto}[5]
\bkvilbox{冷却フィンの温度：両はしの条件と $\cosh$}
パソコンの CPU の熱を逃がすフィン（うすい金属の板）を考えます。根もとからの距離を $x$〔m〕、そこでの板の温度と部屋の温度の差を $\theta$〔K〕とします。板の中を伝わる熱と、表面から空気へ逃げる熱のつり合いを考えると、
\[ \frac{\dd^2\theta}{\dd x^2}=\beta^2\theta \]
という、$y''=a^2y$ と同じ形の式が成り立ちます（$\beta$（ベータ）〔1/m〕は、板の材料・厚さ・風の強さで決まる数。くわしくは物理『熱』）。条件は、根もとの温度差 $\theta(0)=\theta_{\mathrm b}$ と、先たん（$x=L$）からは熱がほとんど逃げないこと $\theta'(L)=0$ です。6-3 と同じ「両はしの条件」です。答えは
\[ \theta=\theta_{\mathrm b}\,\frac{\cosh\bigl(\beta(L-x)\bigr)}{\cosh\beta L} \]
です。確かめます。
\begin{enumerate}[label={\bkcircled[\bkTF]{\arabic*}}, leftmargin=2.4\zw, itemsep=1pt, topsep=2pt]
\item 2回微分すると、中の $\beta(L-x)$ の微分 $-\beta$ が2回かかって $(-\beta)^2=\beta^2$ 倍。式をみたす。
\item $x=0$ では、分子と分母が同じ $\cosh\beta L$ なので $\theta=\theta_{\mathrm b}$。
\item 1回微分すると $-\beta\theta_{\mathrm b}\dfrac{\sinh\bigl(\beta(L-x)\bigr)}{\cosh\beta L}$。$x=L$ では $\sinh0=0$ なので傾きは0。
\end{enumerate}
\noindent\begin{minipage}[c]{0.56\linewidth}
たとえば $\beta=20$〔1/m〕、$L=\qty{0.05}{m}$（$\beta L=1$）、$\theta_{\mathrm b}=\qty{40}{K}$ なら、先たんの温度差は $\dfrac{40}{\cosh1}=\dfrac{40}{1.543}\approx\qty{25.9}{K}$。先へ行くほど板は冷えて、熱を逃がす力が弱くなります。フィンを長くしすぎても、あまり冷えなくなるわけです。
\end{minipage}\hfill
\begin{minipage}[c]{0.42\linewidth}
\centering
\begin{tikzpicture}
\begin{axis}[bkplot, scale only axis, width=28mm, height=24mm, xmin=0, xmax=0.056, ymin=0, ymax=46, xtick={0.05}, xticklabels={0.05}, ytick={25.9,40}, yticklabels={25.9,40}, xlabel={$x$（m）}, ylabel={$\theta$（K）},
    yticklabel style={fill=white, inner sep=1pt, outer sep=2.3pt}, xticklabel style={fill=white, inner sep=1pt}]
  \addplot[bkRed, line width=1.6pt, domain=0:0.05, samples=50] {40*cosh(20*(0.05-x))/cosh(1)};
\end{axis}
\end{tikzpicture}
\end{minipage}
\end{motto}

\Needspace{14\baselineskip}
\bksubsection{6}{空気の中の落下：$v=v_{\mathrm t}\tanh\dfrac{gt}{v_{\mathrm t}}$}

\label{bk:6-11:rakka}6-5 では、速さの2乗に比例する抵抗を受けて落ちるときの式
\bkvilline{$m\dfrac{\dd v}{\dd t}=mg-cv^2$\qquad（下向きが正）}
を、表計算で解きました。この式は、$\tanh$ を使うと、式でも解けます。はじめ止まっていた（$v(0)=0$）ときの答えは、
\bkvilframe{$v=v_{\mathrm t}\tanh\dfrac{gt}{v_{\mathrm t}}$\qquad $\Bigl(v_{\mathrm t}=\sqrt{\dfrac{mg}{c}}\Bigr)$}
です。なぜ $\tanh$ なのか、まず形から考えます。
\begin{itemize}[leftmargin=1.2\zw, itemsep=2pt, topsep=2pt, label={\color{\bkTT}●}]
\item \textbf{はじめ}：$\tanh u$ は原点の近くで $u$ とほとんど同じ（②）。だから、
\[ v\approx v_{\mathrm t}\times\frac{gt}{v_{\mathrm t}}=gt \]
空気の抵抗がないときの自由落下（5-13）と同じ速さで、速くなります。
\item \textbf{やがて}：$\tanh u$ は1に近づく（②）。だから $v$ は $v_{\mathrm t}$（終端速度）に近づきます。
\end{itemize}
「はじめは傾き1でまっすぐ、やがて一定の値で平らになる」という $\tanh$ の形は、落下の速さの変わり方そのものです。

\bkvilhead{方法1：微分して、式に入れて確かめる}
$u=\dfrac{gt}{v_{\mathrm t}}$ とおきます。合成関数の微分（5-6）で、$\dfrac{\dd u}{\dd t}=\dfrac{g}{v_{\mathrm t}}$ をかけます。
\begin{bkhenkei}[0.62]
\frac{\dd v}{\dd t}=v_{\mathrm t}\times\frac{g}{v_{\mathrm t}}\bigl(1-\tanh^2u\bigr) & $(\tanh u)'=1-\tanh^2u$（④） \\
=g\Bigl(1-\frac{v^2}{v_{\mathrm t}^2}\Bigr) & $\tanh u=\dfrac{v}{v_{\mathrm t}}$ を入れる \\
=g-\frac{g}{v_{\mathrm t}^2}\,v^2=g-\frac{c}{m}\,v^2 & $v_{\mathrm t}^2=\dfrac{mg}{c}$ より $\dfrac{g}{v_{\mathrm t}^2}=\dfrac{c}{m}$ \\
\end{bkhenkei}
最後の式は、運動の式 $m\dfrac{\dd v}{\dd t}=mg-cv^2$ の両辺を $m$ でわったものと同じです。$t=0$ では $\tanh0=0$ なので $v(0)=0$。たしかに答えです。

\bkvilhead{方法2：表計算の答えと重ねる}
スカイダイバー（\qty{80}{kg}、$c=\qty{0.26}{N.s^2/m^2}$、$v_{\mathrm t}=\qty{54.9}{m/s}$）で、式の値と、6-2 のオイラー法（$\Delta t=\qty{0.1}{s}$）の値をくらべます。
\begin{bkcenter}
\begin{bktable}{colspec={Q[l,m,wd=36mm] *{5}{Q[c,m,wd=13mm]}}, rowsep=1.5pt, colsep=2pt}
時間 $t$〔s〕 & 1 & 2 & 3 & 5 & 10 \\
式 $v_{\mathrm t}\tanh\dfrac{gt}{v_{\mathrm t}}$〔m/s〕 & 9.70 & 18.8 & 26.9 & 39.1 & 51.9 \\
オイラー法〔m/s〕 & 9.71 & 18.9 & 27.0 & 39.3 & 52.0 \\
抵抗なし $gt$〔m/s〕 & 9.8 & 19.6 & 29.4 & 49.0 & 98.0 \\
\end{bktable}
\end{bkcenter}

\begin{bkfigure}
\begin{tikzpicture}
\begin{axis}[bkplot, scale only axis, width=96mm, height=44mm, xmin=0, xmax=21, ymin=0, ymax=66, xtick={5,10,15,20}, ytick={20,40,54.9}, yticklabels={20,40,54.9}, xlabel={$t$（s）}, ylabel={$v$（m/s）},
    yticklabel style={fill=white, inner sep=1pt, outer sep=2.3pt}, xticklabel style={fill=white, inner sep=1pt}]
  \addplot[bkGray, dashed, line width=0.7pt, domain=0:20.5] {54.91};
  \addplot[bkcurve2, dashed, line width=1.2pt, domain=0:6.5] {9.8*x};
  \addplot[bkcurve, line width=2pt, domain=0:20.5, samples=80] {54.91*tanh(9.8*x/54.91)};
  \addplot[only marks, mark=*, mark size=1.5pt, bkInk] coordinates {(1,9.71) (2,18.86) (3,26.98) (5,39.30) (10,52.02)};
  \node[font=\sffamily\footnotesize, text=bkOrange!80!black, anchor=east, inner sep=1pt, fill=white] at (axis cs:4.2,47) {抵抗なし $gt$};
  \node[font=\sffamily\footnotesize, text=bkBlue, anchor=north west, inner sep=1pt, fill=white] at (axis cs:7,38) {$v_{\mathrm t}\tanh\dfrac{gt}{v_{\mathrm t}}$};
  \node[font=\sffamily\footnotesize, text=bkGray, anchor=south east, inner sep=1pt, fill=white] at (axis cs:20.4,56.5) {終端速度 $v_{\mathrm t}$};
\end{axis}
\end{tikzpicture}
\bkfigcaption{スカイダイバーの速さ。青い曲線が式の答え、黒い点がオイラー法（$\Delta t=\qty{0.1}{s}$）。はじめは $gt$（オレンジの点線）にそい、やがて終端速度に近づく。}
\end{bkfigure}
表計算の点と式の曲線が、ほとんどぴったり重なりました（小さなずれは、オイラー法の $\Delta t$ によるもの。6-2）。$\tanh$ の式があれば、何秒後の速さでも、表を作らずに1回の計算で求められます。

\Needspace{26\baselineskip}
\begin{motto}[5]
\bkvilbox{変数分離で、$\tanh$ の答えを導く}
$v_{\mathrm t}^2=\dfrac{mg}{c}$ より $g=\dfrac{c}{m}v_{\mathrm t}^2$。運動の式を $m$ でわると $\dfrac{\dd v}{\dd t}=\dfrac{c}{m}\bigl(v_{\mathrm t}^2-v^2\bigr)$。6-4 の変数分離で、$v$ を左、$t$ を右に集めます。
\[ \frac{\dd v}{v_{\mathrm t}^2-v^2}=\frac{c}{m}\,\dd t \]
左の分数を、3-8 の恒等式の考えで2つに分けます（右辺を通分すると左辺にもどる）。
\[ \frac{1}{v_{\mathrm t}^2-v^2}=\frac{1}{2v_{\mathrm t}}\Bigl(\frac{1}{v_{\mathrm t}+v}+\frac{1}{v_{\mathrm t}-v}\Bigr) \]
両辺を積分します（5-16 の置換積分で $\int\frac{\dd v}{v_{\mathrm t}-v}=-\ln(v_{\mathrm t}-v)$）。$v(0)=0$ で積分定数は0になり、
\[ \frac{1}{2v_{\mathrm t}}\ln\frac{v_{\mathrm t}+v}{v_{\mathrm t}-v}=\frac{c}{m}\,t \]
$u=\dfrac{gt}{v_{\mathrm t}}$ とおくと、右辺に $2v_{\mathrm t}$ をかけたものは $2u$ です（$\frac{c}{m}=\frac{g}{v_{\mathrm t}^2}$）。だから $\dfrac{v_{\mathrm t}+v}{v_{\mathrm t}-v}=e^{2u}$。これを $v$ について解き、分子と分母を $e^u$ でわると、
\[ v=v_{\mathrm t}\,\frac{e^{2u}-1}{e^{2u}+1}=v_{\mathrm t}\,\frac{e^{u}-e^{-u}}{e^{u}+e^{-u}}=v_{\mathrm t}\tanh u \]
\end{motto}

\Needspace{8\baselineskip}
\bksubsection{7}{\bkhatenaref の答え：くさりの形は $\cosh$}

\label{bk:6-11:kensui}くさりの、いちばん下の点Oから、ある点Pまでの部分を考えます。くさりは止まっているので、この部分にはたらく力はつり合っています（4-6 の力の分解）。くさり1\,m あたりの重さを $w$〔N/m〕、OからPまでのくさりの長さを $s$〔m〕とします。

\noindent\begin{minipage}[c]{0.5\linewidth}
\raggedright
\begin{itemize}[leftmargin=1.2\zw, itemsep=2pt, topsep=0pt, label={\color{\bkTT}●}]
\item Oでは、左のくさりが真横に引く。その力を $H$〔N〕とする。
\item Pでは、右のくさりが、接線の向きに引く。その力を $T$〔N〕、横とのあいだの角を $\phi$ とする。
\item この部分の重さ $ws$〔N〕が、真下にかかる。
\end{itemize}
横の力のつり合い：$T\cos\phi=H$\par
たての力のつり合い：$T\sin\phi=ws$\par
2つの式をわると $\tan\phi=\dfrac{ws}{H}$。$\tan\phi$ は、Pでの接線の傾き $y'$ です（5-3）。
\end{minipage}\hfill
\begin{minipage}[c]{0.46\linewidth}
\centering
\begin{tikzpicture}[x=1mm, y=1mm]
  \draw[bkInk!35, line width=1.6pt, domain=-14:0, samples=30] plot (\x, {20*(cosh(\x/20)-1)});
  \draw[bkInk, line width=2.2pt, domain=0:26, samples=60] plot (\x, {20*(cosh(\x/20)-1)});
  \fill[bkInk] (0,0) circle (1); \fill[bkInk] (26,19.42) circle (1);
  \node[font=\sffamily\footnotesize, anchor=north, inner sep=1pt] at (0,-1.4) {O};
  \node[font=\sffamily\footnotesize, anchor=north west, inner sep=1pt] at (27,18.6) {P};
  \draw[bk arrow, bkBlue] (0,0) -- (-14,0);
  \node[font=\sffamily\footnotesize, text=bkBlue, anchor=north, inner sep=1pt] at (-9,-0.8) {$H$};
  \draw[bk arrow, bkRed] (26,19.42) -- ++(59.5:16);
  \node[font=\sffamily\footnotesize, text=bkRed, anchor=west, inner sep=1pt] at ($(26,19.42)+(59.5:14)+(1,0)$) {$T$};
  \draw[bkGray, dashed, line width=0.5pt] (26,19.42) -- (40,19.42);
  \draw[bkInk, line width=0.5pt] (33,19.42) arc (0:59.5:7);
  \node[font=\sffamily\footnotesize, anchor=west, inner sep=1pt] at (34.6,22.6) {$\phi$};
  \draw[bk arrow, bkOrange] (13,4.38) -- (13,-12);
  \node[font=\sffamily\footnotesize, text=bkOrange!80!black, anchor=west, inner sep=1pt] at (14,-8) {重さ $ws$};
\end{tikzpicture}
\end{minipage}

\noindent そこで $a=\dfrac{H}{w}$〔m〕とおくと、
\bkvilframe{$y'=\dfrac{s}{a}$\qquad（傾き＝くさりの長さ÷$a$）}
$H$ はくさりのどこでも同じ（横の力のつり合い）なので、$a$ は定数です。$H$〔N〕を $w$〔N/m〕でわるので、$a$ の単位は m です。

\bkvilhead{くさりの長さを、三平方で表す}
\noindent\begin{minipage}[c]{0.62\linewidth}
\raggedright
$y'=\dfrac{s}{a}$ を、もう1回 $x$ で微分すると、$y''=\dfrac{1}{a}\cdot\dfrac{\dd s}{\dd x}$。くさりをとても小さく区切ると、1つ分はほとんどまっすぐで、横 $\dd x$、たて $\dd y$ の直角三角形の斜辺になります（4-4 の三平方）。
\[ \dd s=\sqrt{\dd x^2+\dd y^2}\quad\text{より}\quad\frac{\dd s}{\dd x}=\sqrt{1+y'^2} \]
\end{minipage}\hfill
\begin{minipage}[c]{0.34\linewidth}
\centering
\begin{tikzpicture}[x=1mm, y=1mm]
  \draw[bkInk, line width=1.6pt] (0,0) -- (26,16);
  \draw[bkBlue, line width=1.1pt] (0,0) -- (26,0);
  \draw[bkOrange, line width=1.1pt] (26,0) -- (26,16);
  \draw[bkInk, line width=0.4pt] (23.4,0) -- (23.4,2.6) -- (26,2.6);
  \node[font=\sffamily\footnotesize, text=bkBlue, anchor=north, inner sep=1pt] at (13,-1) {$\dd x$};
  \node[font=\sffamily\footnotesize, text=bkOrange!80!black, anchor=west, inner sep=1pt] at (27,8) {$\dd y$};
  \node[font=\sffamily\footnotesize, anchor=south east, inner sep=1pt] at (12,8.4) {$\dd s$};
\end{tikzpicture}
\end{minipage}\par\vspace{2pt}
（$\dd x$ でわって、$\dfrac{\dd y}{\dd x}=y'$ を使いました。）だから、くさりの形は、次の微分方程式をみたします。
\begin{seikaku}
くさりの小さな1つ分は、本当は少し曲がっているので、まっすぐな斜辺は近似です。区切りを小さくするほど近似は正しくなり、その行き着く先（極限）が曲線の長さ $\displaystyle s=\int_0^x\sqrt{1+y'^2}\,\dd x$ です。ここでの $\dfrac{\dd s}{\dd x}=\sqrt{1+y'^2}$ は、この積分を微分したもの（5-11 の微積分の基本定理）なので、近似ではなく正しい式です。くわしくは大学で学びます。
\end{seikaku}
\bkvilframe{$y''=\dfrac{1}{a}\sqrt{1+y'^2}$}

\bkvilhead{答えは $y=a\Bigl(\cosh\dfrac{x}{a}-1\Bigr)$}
いちばん下の点Oを原点にすると、答えは $y=a\Bigl(\cosh\dfrac{x}{a}-1\Bigr)$ です。式に入れて確かめます。
\begin{enumerate}[label={\bkcircled[\bkTF]{\arabic*}}, leftmargin=2.4\zw, itemsep=2pt, topsep=2pt]
\item 微分する（④）：$y'=a\times\dfrac{1}{a}\sinh\dfrac{x}{a}=\sinh\dfrac{x}{a}$。もう1回：$y''=\dfrac{1}{a}\cosh\dfrac{x}{a}$。
\item 右辺：③の $\cosh^2-\sinh^2=1$ から $1+\sinh^2\dfrac{x}{a}=\cosh^2\dfrac{x}{a}$。$\cosh$ は正なので $\sqrt{1+y'^2}=\cosh\dfrac{x}{a}$。
\item 右辺は $\dfrac{1}{a}\cosh\dfrac{x}{a}$ で、①の $y''$ と同じ。
\item $y(0)=a(1-1)=0$、$y'(0)=\sinh0=0$。Oが底（傾き0）になっている。
\end{enumerate}
くさりの形は、$\cosh$ のグラフを、たても横も $a$ 倍に広げた形です。\bkhatenaref の答えは、「$\cosh$ の曲線」でした。この曲線を\textbf{懸垂線}（けんすいせん）または\textbf{カテナリー}といいます。\bkidx{けんすいせん}{懸垂線}\bkidx{かてなりー}{カテナリー}1691年に、ライプニッツ、ホイヘンス、ヨハン・ベルヌーイが、それぞれこの答えにたどり着きました。カテナリーは、ラテン語で「くさり」を表す catena から来たことばです。

$y'=\dfrac{s}{a}$ から、Oからのくさりの長さも分かります：$s=ay'=a\sinh\dfrac{x}{a}$。

\bkvilhead{やってみようの答え}
はば20\,cm（$x=\pm10$）、たるみ15\,cm なので、次の式をみたす $a$ を表計算でさがします。
\bkvilline{$a\Bigl(\cosh\dfrac{10}{a}-1\Bigr)=15$}
すると $a\approx4.76$\,cm。$x=6$ では、
\bkvilline{$4.76\times\Bigl(\cosh\dfrac{6}{4.76}-1\Bigr)\approx4.3$\,cm}
で、はかった値と合います。

\bkvilhead{6-3 のつり橋（放物線）とのちがい}
6-3 のつり橋では、ケーブルにかかる橋の板の重さが、\textbf{横の長さ $x$} に比例しました。上と同じように考えると $y'=\dfrac{x}{a}$ で、積分すると放物線 $y=\dfrac{x^2}{2a}$ です。くさりでは、重さが\textbf{くさりの長さ $s$} に比例します。はしのほうでは、くさりがななめになって、横の長さより $s$ のほうが長くなるので、重さが多くかかり、はしで急に立ち上がるのです。
\begin{bkfigure}
\begin{tikzpicture}
\begin{axis}[bkplot, scale only axis, width=100mm, height=38mm, xmin=-55, xmax=55, ymin=0, ymax=31, xtick={-50,-25,25,50}, ytick={10,20,27.15}, yticklabels={10,20,27.2}, xlabel={$x$（m）}, ylabel={$y$（m）},
    yticklabel style={fill=white, inner sep=1pt, outer sep=2.3pt}, xticklabel style={fill=white, inner sep=1pt}]
  \addplot[bkcurve2, dashed, line width=1.4pt, domain=-50:50, samples=80] {27.154*(x/50)^2};
  \addplot[bkcurve, line width=2pt, domain=-50:50, samples=100] {50*(cosh(x/50)-1)};
  \node[font=\sffamily\footnotesize, text=bkBlue, anchor=north west, inner sep=1pt, fill=white] at (axis cs:30,8.5) {懸垂線（くさり）};
  \node[font=\sffamily\footnotesize, text=bkOrange!80!black, anchor=south east, inner sep=1pt, fill=white] at (axis cs:-27,12.5) {放物線};
\end{axis}
\end{tikzpicture}
\bkfigcaption{$a=\qty{50}{m}$、はば\qty{100}{m}の懸垂線（たるみ\qty{27.2}{m}）と、同じはば・たるみの放物線。$x=25$ では、6.38\,m と 6.79\,m。}
\end{bkfigure}

\bkvilhead{たるみが小さければ、ほとんど放物線}
5-17 のテイラー展開で、$e^u$ と $e^{-u}$ の展開をたして2でわると、奇数乗の項は消えて、
\bkvilline{$\cosh u=1+\dfrac{u^2}{2!}+\dfrac{u^4}{4!}+\cdots$}
$u$ が小さければ、$\cosh u\approx1+\dfrac{u^2}{2}$。だから
\bkvilline{$y=a\Bigl(\cosh\dfrac{x}{a}-1\Bigr)\approx a\times\dfrac{1}{2}\Bigl(\dfrac{x}{a}\Bigr)^2=\dfrac{x^2}{2a}$}
これは、つり橋の放物線と同じ式です。電線のように、はばにくらべてたるみが小さいとき（$\dfrac{x}{a}$ が小さいとき）は、懸垂線と放物線はほとんど同じです。たとえば、はば\qty{100}{m}で $a=\qty{625}{m}$ の電線では、たるみは懸垂線で2.001\,m、放物線で2.000\,m。ちがいは1\,mm です。ガリレオの「放物線に近い」は、たるみが小さければ正しかったのです。

\begin{motto}[5]
\bkvilbox{$\cos(jx)=\cosh x$}
6-10 の $\cos\theta=\dfrac{e^{j\theta}+e^{-j\theta}}{2}$ の $\theta$ に、虚数 $jx$ を入れてみます。$j\times j=-1$ なので、$e^{j\cdot jx}=e^{-x}$、$e^{-j\cdot jx}=e^{x}$ です。
\[ \cos(jx)=\frac{e^{-x}+e^{x}}{2}=\cosh x \]
「$\cos$ に虚数を入れると $\cosh$ になる」のです。円の関数と双曲線の関数は、複素数の世界ではつながっています。また、$\sinh x$ の逆関数（$\sinh x=y$ から $x$ を求める関数）は $\sinh^{-1}y$ と書き、電卓の hyp キーのあと $\sin^{-1}$ キーで計算できます。
\end{motto}

\begin{senmonsho}[大学の教科書ふうの書き方]
\begin{bkhonbun}
懸垂線の方程式は $y=a\cosh\dfrac{x}{a}$ で与えられる。ここで $a=\dfrac{H}{w}$ であり、$H$ は最下点における張力、$w$ は単位長さあたりの重量である。
\end{bkhonbun}
\bkyaku{大学の本では、いちばん下の点の高さを $a$ にして（$x=0$ で $y=a\cosh0=a$）、$-1$ をつけない形で書くことが多い。この章の $y=a\bigl(\cosh\frac{x}{a}-1\bigr)$ を、たてに $a$ だけずらしたもの。「張力」は、ひもやくさりが引く力。「単位長さあたりの重量」は、1\,m あたりの重さ $w$。}
\end{senmonsho}

\begin{taisetsu}
\def\tsl#1{\makebox[6.4\zw][l]{\sffamily\bfseries\small\color{\bkTT}#1}}
\tsl{定義}$\cosh x=\dfrac{e^{x}+e^{-x}}{2}$、\quad $\sinh x=\dfrac{e^{x}-e^{-x}}{2}$\\[3pt]
\tsl{}$\tanh x=\dfrac{\sinh x}{\cosh x}$\\[3pt]
\tsl{2乗の差}$\cosh^2x-\sinh^2x=1$\\[3pt]
\tsl{}点 $(\cosh t,\ \sinh t)$ は双曲線 $x^2-y^2=1$ の上\\[3pt]
\tsl{微分}$(\cosh x)'=\sinh x$、\quad $(\sinh x)'=\cosh x$\\[3pt]
\tsl{}$(\tanh x)'=1-\tanh^2x$\\[3pt]
\tsl{使い方}$y''=a^2y$ の答えは $C\cosh at+D\sinh at$\\[3pt]
\tsl{}$C=y(0)$、\quad $D=\dfrac{y'(0)}{a}$
\end{taisetsu}

\begin{kotoba}
\small
\bkword{双曲線関数}{$\cosh x$・$\sinh x$・$\tanh x$ のこと。$e^x$ と $e^{-x}$ から作る。}
\bkword{双曲線}{$x^2-y^2=1$ のような曲線。反比例のグラフも双曲線。}
\bkword{懸垂線}{カテナリーともいう。くさりや電線を両はしでつるしたときの形。$y=a\bigl(\cosh\frac{x}{a}-1\bigr)$。}
\end{kotoba}

\bkkigocards{\bkkigocard{$\cosh x$}{ハイパボリック・コサイン}{$\dfrac{e^x+e^{-x}}{2}$。$\sinh x$（ハイパボリック・サイン）は $\dfrac{e^x-e^{-x}}{2}$}{hyperbolic cosine}{6-11}}
{\bkkigocard{$\tanh x$}{ハイパボリック・タンジェント（タンジェント・エイチ）}{$\dfrac{\sinh x}{\cosh x}$。$-1$ と1のあいだ}{hyperbolic tangent}{6-11}}

%% ---------------- 例題 ----------------
\bksection{reidai}{例題でたしかめよう}

\begin{reidai}[$\cosh$・$\sinh$・$\tanh$ の値]
$e=2.718$、$e^{-1}=0.368$ として、$\cosh1$、$\sinh1$、$\tanh1$ を求めましょう。\par
\kangaekata 定義（②）に入れる。$\tanh$ は $\sinh$ を $\cosh$ でわる。\par
\shiki\hangindent=2.6\zw\hangafter=1 $\cosh1=\dfrac{2.718+0.368}{2}=\dfrac{3.086}{2}=1.543$\par
\bkvila $\sinh1=\dfrac{2.718-0.368}{2}=\dfrac{2.350}{2}=1.175$\par
\bkvila $\tanh1=\dfrac{1.175}{1.543}=0.7615\cdots$\par
\kotae\hangindent=2.6\zw\hangafter=1 $\cosh1=1.543$、$\sinh1=1.175$、$\tanh1\approx0.762$
\end{reidai}

\begin{reidai}[$\cosh^2x-\sinh^2x=1$ を使う]
$x>0$ で $\sinh x=\dfrac{3}{4}$ のとき、$\cosh x$ と $\tanh x$ を求めましょう。\par
\kangaekata $\cosh^2x=1+\sinh^2x$（③）。$\cosh x$ はいつも正（①）なので、正のほうをとる。\par
\shiki\hangindent=2.6\zw\hangafter=1 $\cosh^2x=1+\Bigl(\dfrac{3}{4}\Bigr)^2=1+\dfrac{9}{16}=\dfrac{25}{16}$\par
\bkvila $\cosh x>0$ なので $\cosh x=\sqrt{\dfrac{25}{16}}=\dfrac{5}{4}$\par
\bkvila $\tanh x=\dfrac{3}{4}\div\dfrac{5}{4}=\dfrac{3}{4}\times\dfrac{4}{5}=\dfrac{3}{5}$\par
\kotae\hangindent=2.6\zw\hangafter=1 $\cosh x=\dfrac{5}{4}$、$\tanh x=\dfrac{3}{5}$
\end{reidai}

\noindent\begin{minipage}[t]{0.488\linewidth}
\begin{jibun}
$e^2=7.389$、$e^{-2}=0.135$ として、$\cosh2$、$\sinh2$、$\tanh2$ を求めましょう。
\bkjans{$\cosh2=\dfrac{7.389+0.135}{2}=3.762$、$\sinh2=\dfrac{7.389-0.135}{2}=3.627$、$\tanh2=\dfrac{3.627}{3.762}\approx0.964$。}
\end{jibun}
\end{minipage}\hfill
\begin{minipage}[t]{0.488\linewidth}
\begin{jibun}
$x>0$ で $\sinh x=\dfrac{12}{5}$ のとき、$\cosh x$ と $\tanh x$ を求めましょう。
\bkjans{$\cosh^2x=1+\dfrac{144}{25}=\dfrac{169}{25}$ より $\cosh x=\dfrac{13}{5}$。$\tanh x=\dfrac{12}{13}$。}
\end{jibun}
\end{minipage}\par

\begin{reidai}[双曲線関数の微分]
次の関数を微分しましょう。\par
（1）$\cosh3x$ を1回\quad（2）$\sinh2x$ を2回\quad（3）$5\tanh2t$ を1回\par
\kangaekata ④の表と、合成関数の微分（中の微分をかける）。$\tanh$ は $(\tanh u)'=1-\tanh^2u$。\par
\shiki\hangindent=4.3\zw\hangafter=1 （1）$(\cosh3x)'=3\sinh3x$\par
\bkvilq （2）$(\sinh2x)'=2\cosh2x$\par
\bkvilb もう1回：$(2\cosh2x)'=2\times2\sinh2x=4\sinh2x$\par
\bkvilq （3）$(5\tanh2t)'=5\times2\,(1-\tanh^22t)$\par
\kotae\hangindent=4.3\zw\hangafter=1 （1）$3\sinh3x$\quad（2）$4\sinh2x$\quad（3）$10(1-\tanh^22t)$
\end{reidai}
（1）で、$\cos$ のくせで $-3\sinh3x$ としないように注意しましょう。（2）は「2回微分で自分の $2^2=4$ 倍」になっています。

\begin{reidai}[微分方程式をみたすか確かめる]
$y=\tanh t$ が、$y'=1-y^2$ と $y(0)=0$ をみたすことを確かめましょう。\par
\kangaekata 左辺 $y'$ を計算し、右辺に $y=\tanh t$ を入れたものとくらべる。\par
\shiki\hangindent=2.6\zw\hangafter=1 左辺：$y'=(\tanh t)'=1-\tanh^2t$\par
\bkvila 右辺：$1-y^2=1-\tanh^2t$。左辺と同じ\par
\bkvila $y(0)=\tanh0=\dfrac{\sinh0}{\cosh0}=\dfrac{0}{1}=0$\par
\par\vspace{3pt}{\centering
\begin{tikzpicture}
\begin{axis}[bkplot, scale only axis, width=62mm, height=24mm, xmin=-2.6, xmax=2.6, ymin=-1.2, ymax=1.25, xtick={-2,-1,1,2}, ytick={-1,1}, xlabel={$t$}, ylabel={$y$},
    yticklabel style={fill=white, inner sep=1pt, outer sep=2.3pt}, xticklabel style={fill=white, inner sep=1pt}]
  \addplot[bkGray, dashed, line width=0.6pt, domain=-2.6:2.6] {1};
  \addplot[bkGray, dashed, line width=0.6pt, domain=-2.6:2.6] {-1};
  \addplot[bkPurple, line width=1.6pt, domain=-2.5:2.5, samples=60] {tanh(x)};
  \addplot[bkOrange, line width=1pt, domain=1.3:2.3] {tanh(1.8)+(1-tanh(1.8)^2)*(x-1.8)};
  \node[font=\sffamily\footnotesize, text=bkOrange!80!black, anchor=north west, inner sep=1pt, fill=white] at (axis cs:1.0,0.75) {傾き $1-y^2$（小さい）};
  \node[font=\sffamily\footnotesize, text=bkPurple, anchor=north west, inner sep=1pt, fill=white] at (axis cs:-2.4,0.9) {$y=\tanh t$};
\end{axis}
\end{tikzpicture}\par}\vspace{2pt}
\kotae\hangindent=2.6\zw\hangafter=1 どちらもみたす
\end{reidai}

\noindent\begin{minipage}[t]{0.488\linewidth}
\begin{jibun}
次の関数を微分しましょう。\par
（1）$\cosh2x$ を1回\par（2）$\sinh3x$ を2回\par（3）$4\tanh0.5t$ を1回
\bkjans{（1）\nobreak $2\sinh2x$。（2）\nobreak $9\sinh3x$。（3）\nobreak $4\times0.5(1-\tanh^20.5t)=2(1-\tanh^20.5t)$。}
\end{jibun}
\end{minipage}\hfill
\begin{minipage}[t]{0.488\linewidth}
\begin{jibun}
$y=2\tanh3t$ が、$y'=6-1.5y^2$ と $y(0)=0$ をみたすことを確かめましょう。
\bkjans{$y'=6(1-\tanh^23t)=6-6\tanh^23t$。右辺は $6-1.5\times4\tanh^23t=6-6\tanh^23t$ で同じ。$y(0)=2\tanh0=0$。}
\end{jibun}
\end{minipage}\par

\begin{reidai}[$y''=a^2y$ を $\cosh$・$\sinh$ で解く]
（1）$y''=4y$、$y(0)=2$、$y'(0)=6$ をみたす $y$ を、$C\cosh2t+D\sinh2t$ の形で求めましょう。\par
（2）$y=3e^{2t}+e^{-2t}$ を、$C\cosh2t+D\sinh2t$ の形に書きなおしましょう。\par
\kangaekata（1）$C=y(0)$、$D=\dfrac{y'(0)}{a}$（⑤）。（2）$C=A+B$、$D=A-B$（⑤の行き来）。\par
\shiki\hangindent=4.3\zw\hangafter=1 （1）$a^2=4$ なので $a=2$\par
\bkvilb $C=y(0)=2$、$D=\dfrac{y'(0)}{2}=\dfrac{6}{2}=3$\par
\bkvilb 確かめ：$y'=4\sinh2t+6\cosh2t$ で、$y'(0)=0+6=6$\par
\bkvilq （2）$A=3$、$B=1$ なので $C=3+1=4$、$D=3-1=2$\par
\bkvilb 確かめ：$4\cosh2t=2e^{2t}+2e^{-2t}$、$2\sinh2t=e^{2t}-e^{-2t}$。たすと $3e^{2t}+e^{-2t}$\par
\kotae\hangindent=4.3\zw\hangafter=1 （1）$y=2\cosh2t+3\sinh2t$\quad（2）$y=4\cosh2t+2\sinh2t$
\end{reidai}

\begin{jibun}
（1）$y''=9y$、$y(0)=1$、$y'(0)=6$ をみたす $y$ を、$C\cosh3t+D\sinh3t$ の形で求めましょう。\par
（2）$y=5e^{3t}+e^{-3t}$ を、$C\cosh3t+D\sinh3t$ の形に書きなおしましょう。
\bkjans{（1）\nobreak $C=1$、$D=\dfrac{6}{3}=2$。$y=\cosh3t+2\sinh3t$。（2）\nobreak $C=5+1=6$、$D=5-1=4$。$y=6\cosh3t+4\sinh3t$。}
\end{jibun}

\begin{reidai}[スカイダイバーの速さ]
スカイダイバー（$v_{\mathrm t}=\qty{54.9}{m/s}$）が、止まった状態から落ち始めます。$v=v_{\mathrm t}\tanh\dfrac{gt}{v_{\mathrm t}}$ を使って、（1）3秒後の速さ、（2）終端速度の90\,\%になる時刻を求めましょう（$g=\qty{9.8}{m/s^2}$、$\tanh0.535=0.490$、$\tanh1.47=0.900$）。\par
\kangaekata（1）$u=\dfrac{gt}{v_{\mathrm t}}$ を計算して $\tanh u$ を使う。（2）$\tanh u=0.9$ となる $u$ から $t=\dfrac{uv_{\mathrm t}}{g}$。\par
\shiki\hangindent=4.3\zw\hangafter=1 （1）$u=\dfrac{9.8\times3}{54.9}=\dfrac{29.4}{54.9}=0.535$\par
\bkvilb $v=54.9\times\tanh0.535=54.9\times0.490=26.9$\par
\bkvilq （2）$\dfrac{v}{v_{\mathrm t}}=\tanh u=0.9$ なので $u=1.47$\par
\bkvilb $\dfrac{gt}{v_{\mathrm t}}=1.47$ より $t=\dfrac{1.47\times54.9}{9.8}=\dfrac{80.7}{9.8}=8.2$\par
\par\vspace{3pt}{\centering
\begin{tikzpicture}
\begin{axis}[bkplot, scale only axis, width=66mm, height=26mm, xmin=0, xmax=13, ymin=0, ymax=62, xtick={3,8.2}, xticklabels={3,8.2}, ytick={26.9,49.4}, yticklabels={26.9,49.4}, xlabel={$t$（s）}, ylabel={$v$（m/s）},
    yticklabel style={fill=white, inner sep=1pt, outer sep=2.3pt}, xticklabel style={fill=white, inner sep=1pt}]
  \addplot[bkGray, dashed, line width=0.6pt, domain=0:13] {54.91};
  \addplot[bkBlue, line width=1.6pt, domain=0:12.8, samples=60] {54.91*tanh(9.8*x/54.91)};
  \addplot[only marks, mark=*, mark size=1.6pt, bkOrange] coordinates {(3,26.88) (8.25,49.42)};
  \node[font=\sffamily\footnotesize, text=bkGray, anchor=south east, inner sep=1pt] at (axis cs:12.9,55.6) {$v_{\mathrm t}=54.9$};
\end{axis}
\end{tikzpicture}\par}\vspace{2pt}
\kotae\hangindent=4.3\zw\hangafter=1 （1）約\qty{26.9}{m/s}\quad（2）約8.2秒後
\end{reidai}

\begin{jibun}
同じスカイダイバーの、（1）5秒後の速さ、（2）終端速度の95\,\%になる時刻を求めましょう（$\tanh0.892=0.713$、$\tanh1.83=0.950$）。
\bkjans{（1）\nobreak $u=\dfrac{9.8\times5}{54.9}=0.892$。$54.9\times0.713\approx\qty{39.1}{m/s}$。（2）\nobreak $u=1.83$。$t=\dfrac{1.83\times54.9}{9.8}\approx10.3$\,秒後。}
\end{jibun}

\begin{reidai}[懸垂線のたるみと長さ]
$a=\qty{50}{m}$ の懸垂線 $y=a\Bigl(\cosh\dfrac{x}{a}-1\Bigr)$ で、はば\qty{100}{m}（$x=-50$ から50）のときの、たるみと、くさりの長さを求めましょう（$\cosh1=1.543$、$\sinh1=1.175$）。\par
\kangaekata たるみは、はし（$x=50$）の高さ $y$。長さは、Oから片がわが $s=a\sinh\dfrac{x}{a}$（⑦）なので、その2倍。\par
\shiki\hangindent=2.6\zw\hangafter=1 $\dfrac{x}{a}=\dfrac{50}{50}=1$\par
\bkvila たるみ：$50\times(\cosh1-1)=50\times0.543=27.15$\par
\bkvila 長さ：$2\times50\times\sinh1=100\times1.175=117.5$\par
\par\vspace{3pt}{\centering
\begin{tikzpicture}[x=0.6mm, y=0.6mm]
  \draw[bkInk, line width=1.5pt, domain=-50:50, samples=80] plot (\x, {50*(cosh(\x/50)-1)});
  \fill[bkRed] (-50,27.15) circle (1.3); \fill[bkRed] (50,27.15) circle (1.3);
  \draw[bkGray, line width=0.5pt, {Stealth[length=3pt]}-{Stealth[length=3pt]}] (-50,-5) -- (50,-5);
  \node[font=\sffamily\footnotesize, anchor=north, inner sep=1pt] at (0,-6) {はば 100\,m};
  \draw[bkGray, dashed, line width=0.5pt] (-50,27.15) -- (60,27.15);
  \draw[bkGray, line width=0.5pt, {Stealth[length=3pt]}-{Stealth[length=3pt]}] (58,0) -- (58,27.15);
  \draw[bkGray, dashed, line width=0.5pt] (0,0) -- (60,0);
  \node[font=\sffamily\footnotesize, anchor=west, inner sep=1pt] at (59.5,13.6) {たるみ};
\end{tikzpicture}\par}\vspace{2pt}
\kotae\hangindent=2.6\zw\hangafter=1 たるみ 約\qty{27.2}{m}、長さ 約\qty{117.5}{m}
\end{reidai}

\begin{jibun}
$a=\qty{100}{m}$、はば\qty{100}{m}の懸垂線の、たるみと長さを求めましょう（$\cosh0.5=1.128$、$\sinh0.5=0.521$）。
\bkjans{$\dfrac{x}{a}=0.5$。たるみ $100\times0.128=\qty{12.8}{m}$。長さ $2\times100\times0.521=\qty{104.2}{m}$。}
\end{jibun}

\begin{reidai}[電線のたるみ：放物線で近似する]
1\,m あたりの重さが $w=\qty{10}{N/m}$ の電線を、いちばん下での張力 $H=\qty{5000}{N}$ で、はば\qty{50}{m}にはりました。$a$ を求め、たるみを $\dfrac{x^2}{2a}$（⑦の近似）で求めましょう。\par
\kangaekata $a=\dfrac{H}{w}$（⑦）。たるみは、はし $x=25$ での高さ。$\dfrac{x}{a}$ が小さいので放物線で近似できる。\par
\shiki\hangindent=2.6\zw\hangafter=1 $a=\dfrac{5000}{10}=500$〔m〕。$\dfrac{x}{a}=\dfrac{25}{500}=0.05$（小さい）\par
\bkvila たるみ $\approx\dfrac{25^2}{2\times500}=\dfrac{625}{1000}=0.625$\par
\kotae\hangindent=2.6\zw\hangafter=1 $a=\qty{500}{m}$、たるみ 約\qty{0.63}{m}（懸垂線の式でも0.625\,m）
\end{reidai}

\begin{jibun}
$w=\qty{8}{N/m}$ の電線を、$H=\qty{3200}{N}$ で、はば\qty{80}{m}にはりました。$a$ とたるみを求めましょう。
\bkjans{$a=\dfrac{3200}{8}=\qty{400}{m}$。$x=40$ で $\dfrac{40^2}{2\times400}=\dfrac{1600}{800}=\qty{2.0}{m}$。}
\end{jibun}

%% ---------------- まちがえやすいところ・物理工学 ----------------
\bksection{tsumazuki}{まちがえやすいところをチェック}

\begin{tsumazuki}{$\cosh$ の微分に $-$ をつける}
\bkng{$(\cosh x)'=-\sinh x$}
\bkok{$(\cosh x)'=\sinh x$}
$\cos$ の微分のくせです。$\cosh x$ は $x>0$ でふえているので、傾きは正のはず。$x=1$ で $-\sinh1=-1.175$ では、へっていることになってしまいます。
\end{tsumazuki}

\begin{tsumazuki}{2乗の「和」が1だと思う}
\bkng{$\cosh^2x+\sinh^2x=1$}
\bkok{$\cosh^2x-\sinh^2x=1$}
$x=1$ で確かめると、$2.381+1.381=3.762$ で1になりません。ひき算で $2.381-1.381=1$ です。$\cos^2+\sin^2=1$ と符号がちがいます。
\end{tsumazuki}

\begin{tsumazuki}{$\tanh$ が $\tan$ のようにくり返すと思う}
\bkng{$\tanh x$ は周期的で、いくらでも大きくなる}
\bkok{$\tanh x$ はくり返さず、いつも $-1$ と1のあいだ}
$\tan$ は、$\cos$ が0になる所で大きくなりますが、$\cosh$ は0にならない（いつも1以上）ので、$\tanh$ は大きくなりません。
\end{tsumazuki}

\begin{tsumazuki}{たるみが大きいのに、放物線で計算する}
\bkng{$a=\qty{50}{m}$、はば\qty{100}{m}でも、たるみ $\approx\dfrac{50^2}{2\times50}=\qty{25}{m}$}
\bkok{$50\times(\cosh1-1)\approx\qty{27.2}{m}$}
放物線の近似は、$\dfrac{x}{a}$ が小さいときだけ使えます。ここでは $\dfrac{x}{a}=1$ で小さくないので、2\,m もずれます。
\end{tsumazuki}

\begin{butsuri}
\bkvilbox{送電線のたるみ}
\noindent\begin{minipage}[c]{0.6\linewidth}
\raggedright
鉄塔のあいだの送電線は、懸垂線の形です。電線は、電気が流れたり夏の日ざしを受けたりして温まると、少しのびます（熱による伸び）。長さ $s$ がのびると、たるみが大きくなり、電線が下がります。設計では、いちばん暑い日でも、電線が木や建物に近づきすぎないように、懸垂線の式でたるみを計算します。
\end{minipage}\hfill
\begin{minipage}[c]{0.36\linewidth}
\centering
\begin{tikzpicture}[x=0.8mm, y=1mm]
  \foreach \xx in {0,32} {\begin{scope}[xshift=\xx mm]
    \draw[bkGray!80!black, line width=0.9pt] (-4,0) -- (0,30) -- (4,0);
    \draw[bkGray!80!black, line width=0.6pt] (-3,7) -- (3,7) (-2,15) -- (2,15) (-5,26) -- (5,26);
  \end{scope}}
  \draw[bkGray, line width=0.6pt] (-6,0) -- (46,0);
  \draw[bkBlue, line width=1.2pt, domain=-20:20, samples=50] plot ({20+\x}, {20+40*(cosh(\x/40)-1)});
  \draw[bkRed, line width=1.2pt, dash pattern=on 2.5pt off 1.2pt, domain=-20:20, samples=50] plot ({20+\x}, {14+26*(cosh(\x/26)-1)});
  \node[font=\sffamily\footnotesize, text=bkBlue, anchor=south, inner sep=1pt] at (20,20.6) {冬};
  \node[font=\sffamily\footnotesize, text=bkRed, anchor=north, inner sep=1pt] at (20,13.4) {夏};
\end{tikzpicture}
\end{minipage}
\end{butsuri}

\begin{butsuri}
\bkvilbox{アーチ：さかさまの懸垂線}
\noindent\begin{minipage}[c]{0.34\linewidth}
\centering
\begin{tikzpicture}[x=0.8mm, y=0.8mm]
  \fill[bkGray!25] (-24,-1.5) rectangle (24,0);
  \draw[bkInk!70, line width=4pt, domain=-20:20, samples=60] plot (\x, {30-12*(cosh(\x/12)-1)});
\end{tikzpicture}
\end{minipage}\hfill
\begin{minipage}[c]{0.62\linewidth}
\raggedright
くさりは、引っぱる力だけで形を保っています。その形を上下さかさまにすると、こんどは、押す力だけで形を保つアーチになります。石やれんがは、引っぱる力には弱く、押す力には強いので、懸垂線に近い形のアーチは、くずれにくいのです。スペインの建築家ガウディは、おもりをつけたひもをつるした模型をさかさまに見て、教会の形を決めました。アメリカのセントルイスにあるゲートウェイ・アーチ（高さもはばも約\qty{192}{m}）も、懸垂線をもとにした形です。
\end{minipage}
\end{butsuri}



\begin{column}{ガリレオとくさり}
ガリレオは、1638年の本『新科学対話』の中で、つるしたくさりの形は放物線に近いと書きました。ガリレオは放物線を、投げた物の通り道として調べていたので、似た形を見つけておどろいたことでしょう。その後、くさりの形を正確に求める問題は、数学者たちの「挑戦状」になりました。1690年にヤコブ・ベルヌーイがこの問題を出し、翌1691年に、弟のヨハン・ベルヌーイ、ライプニッツ、ホイヘンスが答えを発表しました。ヨハンは、兄が解けなかった問題を解いたことを、とても喜んだと伝えられています。この章で、みなさんも同じ答えにたどり着きました。
\end{column}

%% ---------------- 章末問題 ----------------
\bkendproblems
\bkproblemgroup{基本}{4}

\bkproblem $e^{1.5}=4.482$、$e^{-1.5}=0.223$ として、$\cosh1.5$、$\sinh1.5$、$\tanh1.5$ を求めましょう。
\bkpans{$\cosh1.5=\dfrac{4.705}{2}=2.3525$、$\sinh1.5=\dfrac{4.259}{2}=2.1295$、$\tanh1.5=\dfrac{2.1295}{2.3525}\approx0.905$。}

\bkproblem $x>0$ で $\sinh x=\dfrac{5}{12}$ のとき、$\cosh x$ と $\tanh x$ を求めましょう。
\bkpans{$\cosh^2x=1+\dfrac{25}{144}=\dfrac{169}{144}$ より $\cosh x=\dfrac{13}{12}$。$\tanh x=\dfrac{5}{13}$。}

\bkproblem 次の関数を微分しましょう。\newline（1）$\cosh4x$ を1回\quad（2）$\sinh5x$ を2回\quad（3）$3\tanh2t$ を1回
\bkpans{（1）\nobreak $4\sinh4x$。（2）\nobreak $25\sinh5x$。（3）\nobreak $6(1-\tanh^22t)$。}

\bkproblem $y=3\tanh2t$ が、$y'=6-\dfrac{2}{3}y^2$ と $y(0)=0$ をみたすことを確かめましょう。
\bkpans{$y'=6(1-\tanh^22t)=6-6\tanh^22t$。右辺は $6-\dfrac{2}{3}\times9\tanh^22t=6-6\tanh^22t$。$y(0)=0$。}

\bkproblem $y''=16y$、$y(0)=3$、$y'(0)=4$ をみたす $y$ を、$C\cosh4t+D\sinh4t$ の形で求めましょう。
\bkpans{$C=3$、$D=\dfrac{4}{4}=1$。$y=3\cosh4t+\sinh4t$。}

\bkproblem スカイダイバー（$v_{\mathrm t}=\qty{54.9}{m/s}$）の、落ち始めて2秒後の速さを求めましょう（$g=\qty{9.8}{m/s^2}$、$\tanh0.357=0.343$）。
\bkpans{$u=\dfrac{9.8\times2}{54.9}\approx0.357$。$54.9\times0.343\approx\qty{18.8}{m/s}$。}

\bkproblem $a=\qty{20}{m}$、はば\qty{40}{m}の懸垂線の、たるみと長さを求めましょう（$\cosh1=1.543$、$\sinh1=1.175$）。
\bkpans{$\dfrac{x}{a}=1$。たるみ $20\times0.543\approx\qty{10.9}{m}$。長さ $2\times20\times1.175=\qty{47.0}{m}$。}

\bkproblem $e^{0.5}=1.649$、$e^{-0.5}=0.607$ として、$\cosh0.5$、$\sinh0.5$、$\tanh0.5$ を求めましょう。
\bkpans{$\cosh0.5=\dfrac{2.256}{2}=1.128$、$\sinh0.5=\dfrac{1.042}{2}=0.521$、$\tanh0.5=\dfrac{0.521}{1.128}\approx0.462$。}

\bkproblemgroup{チャレンジ}{5}

\bkproblem $y=e^{t}-3e^{-t}$ を、$C\cosh t+D\sinh t$ の形に書きなおしましょう。
\bkpans{$A=1$、$B=-3$。$C=A+B=-2$、$D=A-B=4$。$y=-2\cosh t+4\sinh t$。}

\bkproblem 質量\qty{70}{kg}、$c=\qty{0.25}{N.s^2/m^2}$ の人が、止まった状態から落ち始めます。終端速度 $v_{\mathrm t}$ と、4秒後の速さを求めましょう（$g=\qty{9.8}{m/s^2}$、$\sqrt{2744}=52.4$、$\tanh0.748=0.634$）。
\bkpans{$v_{\mathrm t}=\sqrt{\dfrac{70\times9.8}{0.25}}=\sqrt{2744}\approx\qty{52.4}{m/s}$。$u=\dfrac{9.8\times4}{52.4}\approx0.748$。$52.4\times0.634\approx\qty{33.2}{m/s}$。}

\bkproblem 1\,m あたりの重さ $w=\qty{12}{N/m}$ の電線を、$H=\qty{7200}{N}$ で、はば\qty{60}{m}にはりました。$a$ とたるみを求めましょう。
\bkpans{$a=\dfrac{7200}{12}=\qty{600}{m}$。$x=30$ で $\dfrac{30^2}{2\times600}=\dfrac{900}{1200}=\qty{0.75}{m}$。}

\bkproblem $a=\qty{40}{m}$、はば\qty{40}{m}の懸垂線の、たるみと長さを求めましょう（$\cosh0.5=1.128$、$\sinh0.5=0.521$）。
\bkpans{$\dfrac{x}{a}=\dfrac{20}{40}=0.5$。たるみ $40\times0.128\approx\qty{5.1}{m}$。長さ $2\times40\times0.521\approx\qty{41.7}{m}$。}

\begin{furikaeri}
\begin{checklist}
\item $\cosh x$・$\sinh x$・$\tanh x$ を $e^x$ と $e^{-x}$ で表し、値とグラフの形を説明できる。\bkfurref{bk:6-11:teigi}
\item $\cosh^2x-\sinh^2x=1$ と、微分のきまりを使える。\bkfurref{bk:6-11:bibun}
\item $y''=a^2y$ の答えを $\cosh$・$\sinh$ で書き、落下や電線の形に使える。\bkfurref{bk:6-11:kaku}
\end{checklist}
\end{furikaeri}

\bkfukidashi{mina}{$\cos$ と $\sin$ は「回る・ゆれる」、$\cosh$ と $\sinh$ は「ふえる・たるむ」。次の第3部では、ゆれと $e^{x}$ が組み合わさった動きを、式で解くよ。}
